
\section{Planteamiento del problema}

La empresa distribuidora de biblias y literatura cristiana enfrenta una disminución significativa en sus ventas, situación que se atribuye a factores como el aumento de los costos operativos y la intensificación de la competitividad en el mercado. Este panorama pone de manifiesto la necesidad de evaluar y fortalecer las estrategias que permitan retener a los clientes y maximizar su valor para la organización.

Actualmente, la empresa no cuenta con acceso a información clave sobre el comportamiento de compra de sus clientes, como su intención de compra o el volumen total de transacciones realizadas, datos que usualmente podrían obtenerse mediante encuestas. Sin embargo, estas herramientas representan un costo elevado y están sujetas a sesgos en las respuestas, limitando su efectividad y su alcance. Esta falta de visibilidad en el comportamiento de los clientes dificulta la evaluación precisa de su lealtad y restringe las proyecciones de crecimiento potencial, lo que debilita la capacidad de la empresa para adaptarse a las demandas del mercado \cite{glady2009predicting}.

Además, la ausencia de mecanismos basados en ciencia de datos agrava las limitaciones informativas, reduciendo la capacidad de la distribuidora para identificar patrones clave en el comportamiento de los clientes y formular estrategias competitivas. Si esta problemática persiste, la empresa corre el riesgo de enfrentar serias dificultades financieras que comprometan su sostenibilidad y su capacidad de crecimiento en el mediano y largo plazo.